\textbf{Phần 1. Trắc nghiệm lựa chọn}
\setcounterpageref{ex}{0}

\begin{ex} %[2D8N43]
Kết quả của phép tính $8 : 2$ là
\choice
{$3$}
{$5$}
{$6$}
{\True $4$}
\loigiai{
$8 : 2 = 4$.\\
\textbf{Đáp án của câu hỏi là D}
}
\end{ex}

\begin{ex} %[2D8N43]
Kết quả của phép tính $14 : 2$ là
\choice
{$6$}
{\True $7$}
{$8$}
{$9$}
\loigiai{
$14 : 2 = 7$.\\
\textbf{Đáp án của câu hỏi là B}
}
\end{ex}

\begin{ex} %[2D8H43]
Phép tính nào sau đây có kết quả bằng 9?
\choice
{\True $18 : 2$}
{$16 : 2$}
{$20 : 2$}
{$14 : 2$}
\loigiai{
$18 : 2 = 9$.\\
\textbf{Đáp án của câu hỏi là A}
}
\end{ex}

\begin{ex} %[2D8H43]
Phép tính nào trong bảng chia 2 có kết quả bé nhất?
\choice
{$4 : 2$}
{$6 : 2$}
{\True $2 : 2$}
{$8 : 2$}
\loigiai{
$2 : 2 = 1$ (bé nhất).\\
\textbf{Đáp án của câu hỏi là C}
}
\end{ex}

\begin{ex} %[2D8H43]
Trong bảng chia 2, các số bị chia đều là
\choice
{Các số lẻ}
{Các số tròn chục}
{Các số tận cùng là 5}
{\True Các số chẵn từ 2 đến 20}
\loigiai{
Các số bị chia là các số chẵn: 2, 4, 6, 8, ..., 20.\\
\textbf{Đáp án của câu hỏi là D}
}
\end{ex}

\begin{ex} %[2D8H43]
Điền số thích hợp: $16 : 2 = \dots$
\choice
{\True $8$}
{$7$}
{$9$}
{$6$}
\loigiai{
$16 : 2 = 8$.\\
\textbf{Đáp án của câu hỏi là A}
}
\end{ex}

\begin{ex} %[2D8V43]
Có 12 cái kẹo chia đều cho 2 bạn. Mỗi bạn nhận được số kẹo là
\choice
{$5$ cái}
{$7$ cái}
{\True $6$ cái}
{$8$ cái}
\loigiai{
$12 : 2 = 6$ cái kẹo.\\
\textbf{Đáp án của câu hỏi là C}
}
\end{ex}

\begin{ex} %[2D8V43]
Một sợi dây dài $20\text{ dm}$ cắt làm 2 đoạn bằng nhau. Mỗi đoạn dài là
\choice
{$8\text{ dm}$}
{$12\text{ dm}$}
{$15\text{ dm}$}
{\True $10\text{ dm}$}
\loigiai{
$20 : 2 = 10\text{ dm}$.\\
\textbf{Đáp án của câu hỏi là D}
}
\end{ex}

\begin{ex} %[2D8V43]
Tính giá trị biểu thức: $18 : 2 + 15$.
\choice
{$22$}
{\True $24$}
{$26$}
{$28$}
\loigiai{
$9 + 15 = 24$.\\
\textbf{Đáp án của câu hỏi là B}
}
\end{ex}

\begin{ex} %[2D8V43]
Điền dấu thích hợp: $14 : 2 \dots 16 : 2$.
\choice
{$>$}
{$=$}
{\True $<$}
{Không so sánh được}
\loigiai{
$14 : 2 = 7 < 16 : 2 = 8$.\\
\textbf{Đáp án của câu hỏi là C}
}
\end{ex}

\begin{ex} %[2D8V43]
Có 10 lít nước mắm rót đều vào 2 can. Mỗi can có số lít nước mắm là
\choice
{$4\text{ l}$}
{\True $5\text{ l}$}
{$6\text{ l}$}
{$7\text{ l}$}
\loigiai{
$10 : 2 = 5\text{ l}$.\\
\textbf{Đáp án của câu hỏi là B}
}
\end{ex}

\begin{ex} %[2D8C43]
Có 18 quả cam chia đều vào 2 rổ. Mẹ lấy ở một rổ ra 3 quả. Hỏi rổ đó còn lại bao nhiêu quả cam?
\choice
{\True $6$ quả}
{$5$ quả}
{$7$ quả}
{$8$ quả}
\loigiai{
Mỗi rổ có: $18 : 2 = 9$ quả.\\
Còn lại: $9 - 3 = 6$ quả.\\
\textbf{Đáp án của câu hỏi là A}
}
\end{ex}

\textbf{Phần 2. Trắc nghiệm đúng sai}
\setcounterpageref{ex}{0}

\begin{ex} %[2D8H43]
Xét tính đúng sai của các phép tính trong bảng chia 2:
\choiceTF
{\True $6 : 2 = 3$.}
{\True $12 : 2 = 6$.}
{$18 : 2 = 8$.}
{\True $18 : 2 = 9$.}
\loigiai{
\begin{itemize}
    \item $6 : 2 = 3$: Đúng.
    \item $12 : 2 = 6$: Đúng.
    \item $18 : 2 = 9 \neq 8$. Ý 3 Sai, Ý 4 Đúng.
\end{itemize}
\begin{itemize}
    \item a) \textbf{Đúng}.
    \item b) \textbf{Đúng}.
    \item c) \textbf{Sai}.
    \item d) \textbf{Đúng}.
\end{itemize}
}
\end{ex}

\begin{ex} %[2D8V43]
Có 16 học sinh chia đều thành 2 tổ. Xét tính đúng sai:
\choiceTF
{\True Mỗi tổ có 8 học sinh ($16 : 2 = 8$).}
{\True Số học sinh của hai tổ bằng nhau.}
{Nếu chuyển 2 bạn từ tổ 1 sang tổ 2 thì tổ 2 có 9 bạn.}
{\True Khi đó tổ 2 có $8 + 2 = 10$ bạn.}
\loigiai{
\begin{itemize}
    \item Mỗi tổ 8 bạn: Đúng.
    \item Hai tổ bằng nhau: Đúng.
    \item $8 + 2 = 10 \neq 9$. Ý 3 Sai, Ý 4 Đúng.
\end{itemize}
\begin{itemize}
    \item a) \textbf{Đúng}.
    \item b) \textbf{Đúng}.
    \item c) \textbf{Sai}.
    \item d) \textbf{Đúng}.
\end{itemize}
}
\end{ex}

\textbf{Phần 3. Trả lời ngắn}
\setcounterpageref{ex}{0}

\begin{ex} %[2D8H43]
Tính: $10 : 2$.
\shortans{5}
\loigiai{
\textbf{Đáp số:} $5$\\
$10 : 2 = 5$.
}
\end{ex}

\begin{ex} %[2D8H43]
Tính: $16 : 2$.
\shortans{8}
\loigiai{
\textbf{Đáp số:} $8$\\
$16 : 2 = 8$.
}
\end{ex}

\begin{ex} %[2D8V43]
Tính: $20 : 2$.
\shortans{10}
\loigiai{
\textbf{Đáp số:} $10$\\
$20 : 2 = 10$.
}
\end{ex}

\begin{ex} %[2D8V43]
Tính giá trị biểu thức: $14 : 2 + 18$.
\shortans{25}
\loigiai{
\textbf{Đáp số:} $25$\\
$7 + 18 = 25$.
}
\end{ex}

\begin{ex} %[2D8V43]
Có 12 cái bánh chia đều cho 2 bạn. Hỏi mỗi bạn được bao nhiêu cái bánh?
\shortans{6}
\loigiai{
\textbf{Đáp số:} $6$\\
$12 : 2 = 6$ cái bánh.
}
\end{ex}

\begin{ex} %[2D8C43]
Một đoạn dây dài $18\text{ cm}$ gập đôi lại. Hỏi độ dài mỗi nửa đoạn dây là bao nhiêu xăng-ti-mét?
\shortans{9}
\loigiai{
\textbf{Đáp số:} $9$\\
$18 : 2 = 9\text{ cm}$.
}
\end{ex}
